\documentclass[10pt]{article}
\usepackage[T1]{fontenc}
\usepackage[utf8]{inputenc}
\usepackage[french]{babel}
\usepackage{amsmath,amssymb,mathrsfs,mathpazo}
\usepackage{geometry}
\geometry{a4paper,margin=2cm}
\usepackage{graphicx,array,tabularx,tabularray,colortbl,multirow,multicol,booktabs,enumitem,pifont,ulem,fancybox,diagbox,cancel}
\usepackage{pstricks,pstricks-add,pst-plot,pst-tree,pst-node,pst-text,pst-func,pst-math,pst-bspline,pst-3dplot,pst-eucl,pst-solides3d}
\let\pscalculate\undefined
\usepackage{pst-bezier}
\usepackage{tikz,pgfplots,tkz-tab}
\usetikzlibrary{arrows,arrows.meta,patterns,calc,positioning,trees,decorations.markings,shapes.geometric,shapes.misc,shapes.arrows,intersections,decorations.pathmorphing,decorations.pathreplacing,backgrounds,shadows,babel}
\pgfplotsset{compat=1.18}
\usepgfplotslibrary{fillbetween}
\usepackage[np]{numprint}
\usepackage{listings,makecell,pas-tableur,twoopt}
\let\barre\undefined
\usepackage[pstricks]{tablvar}
\definecolor{darkgreen}{RGB}{0,100,0}
\providecommand{\red}{\color{red}}
\providecommand{\blue}{\color{blue}}
\providecommand{\green}{\color{green}}
\providecommand{\black}{\color{black}}
\providecommand{\white}{\color{white}}
\providecommand{\gray}{\color{gray}}
\providecommand{\cyan}{\color{cyan}}
\providecommand{\magenta}{\color{magenta}}
\providecommand{\yellow}{\color{yellow}}
\providecommand{\violet}{\color{violet}}
\providecommand{\orange}{\color{orange}}
\providecommand{\euro}{\texteuro}
\providecommand{\EUR}{\texteuro}
\providecommand{\ieme}{\textsuperscript{e}}
\providecommand{\degres}{\ensuremath{^\circ}}
\providecommand{\degree}{\ensuremath{^\circ}}
\providecommand{\vv}[1]{\overrightarrow{#1}}
\providecommand{\Coord}[1]{#1}
\providecommand{\decosix}{}
\providecommand{\starredbullet}{\ensuremath{\bullet}}
\pagestyle{empty}
\setlength{\parindent}{0pt}
\usepackage{adjustbox}\DeclareUnicodeCharacter{2236}{:}\DeclareUnicodeCharacter{2190}{\ensuremath{\leftarrow}}
\providecommand{\R}{\mathbb{R}}
\providecommand{\N}{\mathbb{N}}
\providecommand{\Z}{\mathbb{Z}}
\providecommand{\Q}{\mathbb{Q}}
\providecommand{\D}{\mathbb{D}}
\providecommand{\C}{\mathbb{C}}
\providecommand{\e}{{\mathrm{e}}}
\providecommand{\vect}[1]{\overrightarrow{#1}}
\providecommand{\vectt}[1]{\overrightarrow{\mathrm{#1}}}
\providecommand{\vvt}[1]{\overrightarrow{\mathrm{#1}}}
\providecommand{\np}[2][]{{#2~#1}}
\providecommand{\vv}[1]{{\overrightarrow{#1}}}
\renewcommand{\d}{{\mathrm{d}}}
\providecommand{\degres}{{^\circ}}
\providecommand{\degre}{{^\circ}}
\providecommand{\blue}{\color{blue}}
\providecommand{\red}{\color{red}}
\providecommand{\enskip}{\quad}
\providecommand{\llbracket}{\lbrack\!\lbrack}
\providecommand{\rrbracket}{\rbrack\!\rbrack}
\providecommand{\euro}{{\text{€}}}
\providecommand{\no}{{\text{n°}}}
\providecommand{\interval}[2]{{[#1;#2]}}
\providecommand{\intervFO}[2]{{[#1;#2[}}
\providecommand{\dotfill}{\ldots}

\providecommand{\ds}{\displaystyle}
\providecommand{\textding}[1]{#1}
\providecommand{\Oij}{\(\mathrm{O};\vec{\imath},\vec{\jmath}\)}
\providecommand{\Oijk}{\(\mathrm{O};\vec{\imath},\vec{\jmath},\vec{k}\)}
\providecommand{\Ouv}{\(\mathrm{O};\vec{u},\vec{v}\)}
\providecommand{\intervFF}[2]{\ensuremath{[#1;#2]}}
\renewcommand{\euro}{\ifmmode\text{\texteuro}\else\texteuro\fi}
\DeclareUnicodeCharacter{00B0}{\ifmmode^{\circ}\else\textdegree\fi}
\DeclareUnicodeCharacter{20AC}{\ifmmode\text{\texteuro}\else\texteuro\fi}
% Énoncés : Éducation nationale. Transcription et corrigés : APMEP, Denis Vergès et contributeurs. Découpage et classement : CoopMaths, Sébastien Lozano, Éric Elter et contributeurs. Adaptation de la présentation : Maths Explorer.
% CC BY-SA : https://coopmaths.fr/www/about/
% https://coopmaths.fr/alea/static/bac/2025/tex/bac_2025_11_sujet1_caledonie_4_cor.tex
\begin{document}
\noindent\textit{APMEP / CoopMaths. CC BY-SA. Présentation adaptée par Maths Explorer.}\par\medskip
\smallskip

\begin{enumerate}
\item Puisque $x > 0$, on peut écrire $f(x) = x^2\left(\dfrac{\ln (x)}{x^2}  - 1\right)$.

On sait (croissances comparées) que $\displaystyle\lim_{x \to - \infty} \dfrac{\ln (x)}{x^2} = 0$, donc que $\displaystyle\lim_{x \to - \infty} \dfrac{\ln (x)}{x^2} - 1 = - 1$.

Comme $\displaystyle\lim_{x \to - \infty}x^2 = + \infty$, on obtient finalement par produit de limites 

$\displaystyle\lim_{x \to - \infty} x^2 \left(\dfrac{\ln (x)}{x^2}  - 1\right) = - \infty$.

\hfill\textbf{L'affirmation 1 est vraie.}

\item $f(x)=2\,\cos\left (x\right )- \sin\left (x\right )$, donc:
$f'(x)= -2\,\sin\left (x\right ) - \cos\left (x\right )$. Donc:

$\begin{aligned}[t]
- 2f'(x) + 3f(x) 
& = -2\left ( -2\,\sin\left (x\right ) - \cos\left (x\right )\right ) +3 \left (2\,\cos\left (x\right )- \sin\left (x\right ) \right )\\
& = 4\,\sin\left (x\right ) +2\,\cos\left (x\right ) +6\,\cos\left (x\right ) -3\,\sin\left (x\right )\\
 & = \sin\left (x\right ) +8\,\cos\left (x\right )
\end{aligned}$

Donc la fonction $f$ est solution de l'équation différentielle $-2y' +3y =  \sin\left (x\right ) +8\,\cos\left (x\right )$.

\hfill\textbf{L'affirmation 2 est vraie.}

\item On a $u_0 = 25$ et donc $u_1 = \ln (3 \times 25 + 1) = \ln 76 \approx 4,3 < 25$.

On peut donc supposer que la suite est décroissante ce que l'on démontre par récurrence : soit $(P_n) : u_{n+1} < u_n$

\textit{Initialisation} : on a vu que $u_1 < u_0$ donc $P_0$ est vraie.

\textit{Hérédité} : soit $n \in \N$ et supposons que $u_{n+1} < u_n$ : on a succesivement :
\begin{eqnarray*}
u_{n+1} < u_n &\iff & 3u_{n+1} < 3u_n\\
&\iff &3u_{n+1} + 1 < 3u_n + 1\\
&\iff &\ln \left(3u_{n+1} + 1\right) < \ln \left(3u_n + 1\right)\:\text{par croissance de la fonction ln}\\
&\iff &u_{n+2} < u_{n+1}
\end{eqnarray*}
La relation est vraie au rang $n+1$.

\textit{Conclusion} : La relation est vraie au rang 0 et si elle est vraie au rang $n \in \N$, elle l'est aussi au rang $n + 1$ : d'après le principe de récurrence quel que soit $\begin{gathered}n \in \N, \\\: u_{n+1} < u_n\end{gathered}$ : ceci montre que la suite $(u_n)$ est décroissante.

\hfill\textbf{L'affirmation 3 est vraie.}

\item Une application affine est de la forme $h(x) = ax + b$, avec $a \in \R$ et $b \in \R$.

Donc $k(x) = x^4 + x^2 + ax + b$ ; $k$ est une fonction polynôme dérivable sur $\R$ et sur cet intervalle on a successivement :

$k'(x) = 4x^3 + 2x + a$, puis

$k''(x) = 12x^2 + 2$, comme un carré est supérieur ou égal à zéro, il en résulte que $k''(x) \geqslant 2 > 0$.

Sur $\R$, la dérivée seconde est supérieure à zéro : la fonction $k$ est convexe sur $\R$.

\hfill\textbf{L'affirmation 4 est vraie.}

\item Avec 5 lettres différentes le nombre d'anagrammes est égal à $5! = 120$.

Comme EULER possède deux lettres identiques le nombre d'anagrammes est égal à 

$\dfrac{5!}{2!} = \dfrac{120}{2} = 60$.

\hfill\textbf{L'affirmation 5 est fausse.}
\end{enumerate}

\end{document}
